The ElGamal cryptosystem is a public key encryption scheme built from exponentiation in a cyclic group. Its security is tied to the difficulty of distinguishing Diffie--Hellman tuples from random ones. In the presentation used in this chapter, the message is protected by combining a random group-derived pad with the plaintext, so the ciphertext changes each time the same message is encrypted.

\subsection{Key Generation}

Let $G$ be a cyclic group of prime order $q$ generated by an element $g$. The receiver chooses a secret exponent
\[
a \xleftarrow{\$} \{1,2,\dots,q-1\}
\]
and computes
\[
h = g^a.
\]
The public key is
\[
pk = (G,g,h),
\]
and the secret key is
\[
sk = a.
\]
The value $h$ is public, but recovering $a$ from $g$ and $g^a$ is believed to be hard in an appropriate group. This is the basic one-way structure underlying ElGamal.

\subsection{Encryption}

To encrypt a message $m$, the sender chooses a fresh random exponent
\[
r \xleftarrow{\$} \{1,2,\dots,q-1\}.
\]
The ciphertext is then formed as
\[
C = (u,v) = (g^r, h^r \cdot m).
\]
The first component $u=g^r$ reveals the randomizer in exponentiated form, while the second component masks the message by multiplying it with the shared term $h^r = g^{ar}$. The crucial feature is that the value $r$ is newly sampled for every encryption, so even identical plaintexts produce different ciphertexts.

\begin{figure}[h]
\centering
\begin{tikzpicture}[>=Latex]
    \node[draw,rounded corners,fill=gray!10,minimum width=2.8cm,minimum height=0.9cm] (sender) at (0,0) {Sender};
    \node[draw,rounded corners,fill=gray!10,minimum width=3.2cm,minimum height=0.9cm] (pub) at (6.7,0) {Public key $(G,g,h)$};
    \node[draw,rounded corners,fill=gray!10,minimum width=2.9cm,minimum height=0.9cm] (recv) at (13.0,0) {Receiver};
    \draw[->,thick] (sender.east) -- node[above]{choose random $r$} (pub.west);
    \draw[->,thick] (pub.east) -- node[above]{send $(g^r,h^r m)$} (recv.west);
\end{tikzpicture}

\vspace{0.4em}

\fbox{%
\begin{minipage}{0.86\linewidth}
\centering
\textbf{ElGamal intuition:} the sender uses fresh randomness $r$ to derive a one-time masking factor from the receiver's public key, so the same plaintext does not map to the same ciphertext twice.
\end{minipage}%
}
\caption{High-level ElGamal encryption flow.}
\end{figure}

\subsection{Decryption}

Given a ciphertext
\[
(u,v) = (g^r,h^r \cdot m),
\]
the receiver uses the secret key $a$ to compute
\[
u^a = (g^r)^a = g^{ra} = h^r.
\]
The message is then recovered by removing the mask:
\[
m = \frac{v}{u^a}.
\]
Since
\[
\frac{h^r \cdot m}{(g^r)^a} = \frac{g^{ar}\cdot m}{g^{ar}} = m,
\]
decryption is correct.

\subsection{Security Analysis}

The important security claim is that ElGamal is \textbf{semantically secure} under the \textbf{Decisional Diffie--Hellman assumption}. The structure of the ciphertext is
\[
(g^r, h^r \cdot m) = (g^r, g^{ar}\cdot m).
\]
If an adversary could distinguish encryptions of two chosen messages, then that adversary would be able to determine whether a tuple of the form
\[
(g,g^a,g^r,T)
\]
contains
\[
T = g^{ar}
\]
or instead contains a random group element. This gives the reduction idea:
\[
\text{Breaking ElGamal semantic security} \Longrightarrow \text{distinguishing a DDH tuple}.
\]

The intuition is straightforward. If $T=g^{ar}$, then the challenge ciphertext is a valid ElGamal encryption. If $T$ is random, then the masking term behaves like an independent one-time pad in the group setting, so the adversary gains no useful information about the message. Therefore, any non-negligible success in distinguishing messages would contradict DDH hardness.

\subsection{Small Worked Example}

For illustration, consider a small example that is mathematically easy to verify by hand. Let the group be the multiplicative group modulo $23$, and choose generator
\[
g = 5.
\]
Suppose the receiver chooses
\[
a = 6.
\]
Then
\[
h = g^a = 5^6 \equiv 8 \pmod{23},
\]
so the public key is $(23,5,8)$ and the secret key is $6$.

Now let the message be
\[
m = 10,
\]
and let the sender choose random
\[
r = 7.
\]
The sender computes
\[
u = g^r = 5^7 \equiv 17 \pmod{23}
\]
and
\[
h^r = 8^7 \equiv 12 \pmod{23}.
\]
Thus the second ciphertext component is
\[
v = h^r \cdot m \equiv 12 \cdot 10 \equiv 5 \pmod{23}.
\]
The ciphertext is therefore
\[
C = (17,5).
\]

To decrypt, the receiver computes
\[
u^a = 17^6 \equiv 12 \pmod{23}.
\]
The inverse of $12$ modulo $23$ is $2$, because
\[
12 \cdot 2 \equiv 1 \pmod{23}.
\]
Hence
\[
m = v \cdot (u^a)^{-1} \equiv 5 \cdot 2 \equiv 10 \pmod{23},
\]
which recovers the original message.

This example is intentionally small and insecure in practice, but it clearly demonstrates the three core stages of ElGamal: \textbf{key generation}, \textbf{randomized encryption}, and \textbf{mask removal through exponentiation with the secret key}.
